\documentclass[11pt]{article}
\usepackage[utf8]{inputenc}
\usepackage[T1]{fontenc}
\usepackage{lmodern}
\usepackage[margin=1in]{geometry}
\usepackage{amsmath,amssymb,mathtools}
\usepackage{graphicx}
\usepackage{grffile}
\usepackage{float}
\usepackage{hyperref}
\usepackage{parskip}
\usepackage{enumitem}
\usepackage{xurl}
\setlength{\parskip}{0.65em}
\setlength{\parindent}{0pt}
\begin{document}
\title{Daily Research Digest --- 2026-04-16}
\author{}
\date{}
\maketitle
\textbf{Topics:} galaxy clusters, galaxies, gravitational lensing, dark matter.
\section*{Synthesis}
\_Digest generation failed: timed out\_
\newpage
\section*{Paper Summaries and Figures}
\subsection*{1. The Radial and Vertical Structure of Molecular Gas in the Edge-On Galaxy NGC 4565}
\textbf{Focus:} galaxies\\
\textbf{Authors:} Grace Krahm, Adam K. Leroy, Jiayi Sun, Kijeong Yim, Eric W. Koch, Tony Wong, Deanne Fisher, Erik Rosolowsky, Karin Sandstrom, Dyas Utomo, Jesse van de Sande, Marie Martig, Amelia Fraser-McKelvie, Michael R. Hayden\\
\textbf{Published:} 2026-04-15T17:56:30+00:00\\
\textbf{PDF:} \href{https://arxiv.org/pdf/2604.14136v1}{https://arxiv.org/pdf/2604.14136v1}\\
\textbf{Summary:}
We present high-resolution (0.94" $\approx$ 55 pc) ALMA CO(2-1) and 13CO(2-1) observations of the highly inclined (i\textasciitilde{}87.5 deg) galaxy NGC 4565 covering out to galactocentric radius Rgal > $\pm$ 17 kpc. The combination of sensitivity and resolution enables the detection of CO emission well into the HI-dominated outer disk while isolating individual molecular clouds across the full extent of the galaxy. Although often described as an edge-on Milky Way analog, the molecular gas profile of NGC 4565 has a central gap which is more similar to M31. The 13CO/12CO ratio remains consistent at 0.086 $\pm$ 0.009 from Rgal = 5-13 kpc. Based on fits to position-velocity slices, the molecular disk remains thin, with a FWHM scale height of 79.1 $\pm$ 1.6 pc measured from the vertical intensity profile with little evidence for vertical flaring. Molecular clouds in NGC 4565 show sizes, linewidths, and surface densities consistent with those found in similar environments in PHANGS-ALMA galaxies and in M31. We identify a prominent star-forming complex on the ring-an overdensity of molecular gas we term the East Ring Pileup. This feature hosts a compact, multiwavelength-bright region, which we call the Jewel. Effects of galaxy inclination on molecular cloud radius, velocity dispersion, surface density, and virial parameter appear as second-order effects that are strongest in velocity dispersion. At this resolution, GMCs are preferentially aligned with the disk of the galaxy and horizontally elongated by a factor of\textasciitilde{}2.
\subsubsection*{Selected figures}
\begin{figure}[H]
\centering
\includegraphics[width=\linewidth,height=0.42\textheight,keepaspectratio]{/home/kf/research-assistant/data/cache/figures/http-arxiv.org-abs-2604.14136v1/figure\_1.png}
\end{figure}
\newpage
\subsection*{2. Search for CBCs with SSM Components in Data from The First Part of LVK Fourth Observing Run}
\textbf{Focus:} dark matter\\
\textbf{Authors:} Ines Bentara\\
\textbf{Published:} 2026-04-15T17:11:07+00:00\\
\textbf{PDF:} \href{https://arxiv.org/pdf/2604.14095v1}{https://arxiv.org/pdf/2604.14095v1}\\
\textbf{Summary:}
Star evolution models predict the lightest compact objects in the universe to have masses greater than that of the Sun. Nonetheless, alternative scenarios could lead to the formation of sub-solar mass (SSM) compact objects, such as primordial black holes (PBHs). The LIGO-Virgo-KAGRA Collaboration (LVK) has performed a search for gravitational-wave (GW) signals from compact binary coalescences (CBCs) including at least one SSM component during the first part of their fourth observing run (O4a), reporting no statistically significant candidate. This non-detection sets upper limits on the merger rate of such systems, which can be used to constrain PBH formation models and the fraction of dark matter (DM) in PBHs. For PBH binaries forming at late times, the fraction of DM in PBHs is constrained to be <= 1 for masses above 0.9 M\_sun in the case of monochromatic mass functions. In the early-formation scenario, this fraction is limited to $\le 7\%$ at 1 $M_\odot$ and $\le 40\%$ at 0.35 $M_\odot$.
\subsubsection*{Selected figures}
\begin{figure}[H]
\centering
\includegraphics[width=\linewidth,height=0.42\textheight,keepaspectratio]{/home/kf/research-assistant/data/cache/figures/http-arxiv.org-abs-2604.14095v1/figure\_1.png}
\end{figure}
\newpage
\subsection*{3. Radial Distribution of Star Formation and Gas-phase Metallicity in Spiral-Elliptical Galaxy Pairs}
\textbf{Focus:} galaxies\\
\textbf{Authors:} Cailu Shi, Shuai Feng, Shiyin Shen, Linlin Li, Wenyuan Cui, Guozhen Hu\\
\textbf{Published:} 2026-04-15T17:07:30+00:00\\
\textbf{PDF:} \href{https://arxiv.org/pdf/2604.14092v1}{https://arxiv.org/pdf/2604.14092v1}\\
\textbf{Summary:}
Using integral field spectroscopy from SDSS-IV MaNGA, we investigate the radial distributions of star formation rate (SFR) and gas-phase metallicity in spiral galaxies that reside in spiral-elliptical (S+E) pairs. Spirals in S+E pairs show suppressed central star formation and elevated metallicities, whereas spirals in spiral-spiral pairs exhibit centrally enhanced star formation and reduced metallicities. The degree of SFR suppression and metallicity enhancement in S+E pairs depends on the masses of the pair members. Spirals with more massive elliptical companions experience stronger star-formation suppression and larger increases in metallicity, while lower-mass spirals show more pronounced metallicity enhancement. In addition, within S+E systems, galaxies with asymmetric gas velocity fields display enhanced SFR and higher metallicities, whereas those with symmetric velocity fields exhibit clear central suppression. Based on these results, we infer that in S+E pairs, the spiral galaxy experiences suppressed gas accretion once it enters the hot circumgalactic medium of its early-type companion, which leads to the observed decline in star-formation activity. When a close encounter takes place, tidal perturbations can compress the remaining cold gas and trigger enhanced star formation, producing rapid chemical enrichment and the associated increase in metallicity.
\subsubsection*{Selected figures}
\begin{figure}[H]
\centering
\includegraphics[width=\linewidth,height=0.42\textheight,keepaspectratio]{/home/kf/research-assistant/data/cache/figures/http-arxiv.org-abs-2604.14092v1/figure\_1.png}
\end{figure}
\newpage
\subsection*{4. Dark energy, spatial curvature, and star formation efficiency from JWST photometric and spectroscopic high-redshift galaxies}
\textbf{Focus:} galaxies\\
\textbf{Authors:} Leonardo Comini, Sunny Vagnozzi, Abraham Loeb\\
\textbf{Published:} 2026-04-15T13:29:50+00:00\\
\textbf{PDF:} \href{https://arxiv.org/pdf/2604.13866v1}{https://arxiv.org/pdf/2604.13866v1}\\
\textbf{Summary:}
Early observations from the James Webb Space Telescope (JWST) have revealed an overabundance of massive high-redshift galaxies, raising the question of whether this points to new physics beyond $Λ$CDM, or an enhanced formation efficiency of massive stars. We revisit this issue going beyond earlier analyses based on direct comparisons to theoretical bounds at a fixed cosmology, by performing a full Bayesian analysis of the most extreme galaxies in the CEERS imaging and FRESCO spectroscopic samples, jointly constraining cosmological parameters and the baryon-to-star conversion efficiency $ε$. We do so not only within the spatially flat $Λ$CDM model, but also in models where the dark energy equation of state $w$ and/or the spatial curvature parameter $Ω_K$ are allowed to vary, carefully discussing the impact of both $w$ and $Ω_K$ on the cumulative comoving stellar mass density. Within the flat $Λ$CDM model, once cosmological parameters are marginalized over, the CEERS sample provides a weak $2\sigma $ lower limit of $ε\gtrsim 0.07$, compatible with astrophysical expectations. In contrast, the FRESCO sample requires $ε\gtrsim 0.5$ at $2\sigma $, with values $ε\lesssim 0.2$ disfavored at $>5\sigma $. These results do not qualitatively change when we allow $w$ and/or $Ω_K$ to vary, with no evidence for deviations from $w=-1$ or $Ω_K=0$. Our results therefore suggest that the origin of the ``JWST tension'' is unlikely to be cosmological, but lies in the astrophysics of galaxy formation.
\newpage
\subsection*{5. Witnessing the onset of stellar winds in Super-Luminous Supernova Hosts: implications for star-formation-driven outflows in low and high-redshift galaxies}
\textbf{Focus:} galaxies\\
\textbf{Authors:} A. Saldana-Lopez, A. Gkini, M. J. Hayes, R. Lunnan, C. A. Carr, M. Huberty, C. Scarlata, S. Schulze, J. Sollerman\\
\textbf{Published:} 2026-04-15T13:29:45+00:00\\
\textbf{PDF:} \href{https://arxiv.org/pdf/2604.13865v1}{https://arxiv.org/pdf/2604.13865v1}\\
\textbf{Summary:}
Direct observational constraints on the earliest, stellar-wind-dominated phases of galactic outflows remain scarce. We present medium-resolution VLT/X-shooter spectroscopy of six Type I superluminous supernova (SLSN-I) host galaxies at z = 0.15-0.51, exploiting the bright SLSN continua as single, down-the-barrel probes of the host interstellar medium. From nebular emission lines we derive dust-corrected star-formation rates as low as 0.06-0.44 Msun yr$^{-1}$, and gas-phase metallicities in the extremely metal-poor regime (less than nine percent solar). Moreover, all hosts exhibit narrow, blueshifted Mg II 2796, 2803 absorption, indicative of the presence of low-ionization outflows along the line of sight. Voigt modeling of the Mg II absorption yields maximum outflow velocities of $v_{max}$ = 37-104 km s$^{-1}$, placing these galaxies systematically below the empirical $v_{max}$-SFR relations for more evolved galaxies of similar SFR. Given the short lifetimes of the SLSN massive progenitors, we argue that these outflows must originate from preceding stellar wind episodes. Assuming a constant-velocity outflow over 3 Myr and spherical symmetry, we infer wind masses M$_{wind}$ = (0.02-1.0) $\times 10^6$ Msun and mass-outflow rates $\dot{M}_{wind} = 0.01-0.33$ Msun yr$^{-1}$, corresponding to mass-loading factors $η<1$. These results indicate that, during the first few Myr of a burst, stellar winds and radiation pressure alone drive slow and weak outflows in low-mass systems, prior to the onset of dominant supernova feedback. Our work provides one of the first empirical constraints on early feedback phases relevant for high-redshift galaxies, and for time-dependent implementations of stellar feedback in galaxy formation simulations.
\newpage
\subsection*{6. Gravitational emissions and light curves of quasi-periodic orbits in Schwarzschild spacetime embedded in a Dehnen-type dark matter halo}
\textbf{Focus:} dark matter\\
\textbf{Authors:} Shijie Tan, Chunhua Jiang, Dan Li, Shiyang Hu, Chen Deng, Wenbin Lin\\
\textbf{Published:} 2026-04-15T13:06:08+00:00\\
\textbf{PDF:} \href{https://arxiv.org/pdf/2604.13832v1}{https://arxiv.org/pdf/2604.13832v1}\\
\textbf{Summary:}
Timelike orbits in curved spacetimes encode intrinsic information about the background geometry and serve as critical probes for investigating gravitational theories and source distributions. In this study, we investigate strictly closed timelike orbits within a Schwarzschild spacetime embedded in a Dehnen-type dark matter halo. By solving the geodesic equations, we identify various configurations of these closed orbits and simulate their corresponding gravitational waves and electromagnetic light curves. Our findings reveal that the morphology of closed orbits is primarily governed by the ratio of the azimuthal period to the radial period. Notably, dark matter halo parameters such as the core scale and density parameters exert a significant amplification effect on the orbital scale, which further induces a discernible phase lag in the gravitational wave signals. Furthermore, although certain orbital structures including the number of leaves remain challenging to distinguish via gravitational wave signals alone, they exhibit identifiable signatures in the characteristic peaks of light curves. These findings reveal the multi-messenger potential of closed orbits in bridging black hole environments and dark matter properties, providing theoretical guidance for future dark matter searches.
\subsubsection*{Selected figures}
\begin{figure}[H]
\centering
\includegraphics[width=\linewidth,height=0.42\textheight,keepaspectratio]{/home/kf/research-assistant/data/cache/figures/http-arxiv.org-abs-2604.13832v1/figure\_1.png}
\end{figure}
\newpage
\subsection*{7. The DECam MAGIC Survey: Investigating the Jet Stellar Stream with Photometric Metallicities}
\textbf{Focus:} dark matter\\
\textbf{Authors:} H Q. Do, A. Chiti, P. S. Ferguson, A. P. Ji, G. Limberg, K. Atzberger, J. Carlin, W. Cerny, A. Drlica-Wagner, G. F. Lewis, T. S. Li, S. L. Martell, G. E. Medina, N. E. D. Noel, A. B. Pace, V. Placco, A. H. Riley, D. J. Sand, G. S. Stringfellow, J. A. Carballo-Bello, D. Erkal, S. E. Koposov, B. Mutlu-Pakdil, M. Navabi, N. Shipp, A. K. Vivas, A. Zenteno, D. Zucker\\
\textbf{Published:} 2026-04-15T00:45:25+00:00\\
\textbf{PDF:} \href{https://arxiv.org/pdf/2604.13374v1}{https://arxiv.org/pdf/2604.13374v1}\\
\textbf{Summary:}
Stellar streams are dynamically fragile structures formed by the tidal disruption of dwarf galaxies and stellar clusters. These objects are valuable tracers of the gravitational potential and accretion history of the Milky Way, and are key probes for the presence and interactions of starless dark matter subhalos. The Jet stream is a $\sim 30^\circ$-long stellar stream that is situated at 30.4 kpc and originates from a disrupted globular cluster.It consists of metal-poor stars that follow a retrograde orbit, reducing the impulse imparted from the Milky Way bar and making it especially sensitive to gravitational perturbations from dark matter subhalos. This paper investigates the known extent of the Jet stream by leveraging photometric metallicities derived from a narrowband filter centered on the Ca II H\&K lines at $\sim$3950A on the Dark Energy Camera (DECam), as part of the Mapping the Ancient Galaxy in CaHK (MAGIC) survey. The wide field-of-view of DECam enables the efficient derivation of photometric metallicities for stars across the full extent of the stream, allowing for a metallicity-based selection to identify likely members. We demonstrate the efficacy of photometric metallicities in isolating stream members when used with Gaia DR3 proper motions, identifying a sample of 213 candidate Jet stream member stars. This then allows for the study of stream morphology, through which we identify a clear fanning of the stream toward the end farther from the Milky Way bar. We provide a list of candidate members, enabling spectroscopic follow-up of the Jet stream to facilitate further studies of its dynamics.
\subsubsection*{Selected figures}
\begin{figure}[H]
\centering
\includegraphics[width=\linewidth,height=0.42\textheight,keepaspectratio]{/home/kf/research-assistant/data/cache/figures/http-arxiv.org-abs-2604.13374v1/figure\_1.png}
\end{figure}
\newpage
\subsection*{8. Observational constraints on nonlocal black holes via gravitational lensing}
\textbf{Focus:} gravitational lensing\\
\textbf{Authors:} Rocco D'Agostino, Vittorio De Falco\\
\textbf{Published:} 2026-04-14T18:48:44+00:00\\
\textbf{PDF:} \href{https://arxiv.org/pdf/2604.13223v1}{https://arxiv.org/pdf/2604.13223v1}\\
\textbf{Summary:}
In this paper, we study the gravitational lensing around the static and spherically symmetric DD black holes, which we recently derived as perturbations of the Schwarzschild geometry within the revised Deser-Woodard theory of nonlocal gravity. We first present general analytical expressions for the deflection angle in both weak- and strong-deflection limits, explicitly relating them to the nonlocal corrections to Schwarzschild spacetime. Subsequently, we analyze lensing observables, such as the post-Newtonian effects and the black hole shadow, to constrain the DD black hole parameter space using current observational bounds. Finally, we perform a joint statistical analysis based on the Fisher information matrix, combining these findings with our previously obtained constraints from quasinormal modes. Our results indicate consistency with general relativity at the $1.13\sigma $ level. This work provides a first assessment of the DD parameter space and offers new insights to probe deviations from Einstein's gravity in view of future larger datasets.
\subsubsection*{Selected figures}
\begin{figure}[H]
\centering
\includegraphics[width=\linewidth,height=0.42\textheight,keepaspectratio]{/home/kf/research-assistant/data/cache/figures/http-arxiv.org-abs-2604.13223v1/figure\_1.png}
\end{figure}
\newpage
\subsection*{9. Searching for axions with quantum interferometry}
\textbf{Focus:} dark matter\\
\textbf{Authors:} Tanmay Kumar Poddar, Michael Spannowsky\\
\textbf{Published:} 2026-04-14T18:03:42+00:00\\
\textbf{PDF:} \href{https://arxiv.org/pdf/2604.13181v1}{https://arxiv.org/pdf/2604.13181v1}\\
\textbf{Summary:}
Quantum phase measurements offer a complementary route to axion searches. We show that axion-photon interactions can imprint both Aharonov-Bohm (AB) and Berry phases in experimentally motivated quantum setups. For a coherently oscillating axion dark matter background, the induced effective current generates a time dependent magnetic flux in an rf-SQUID, leading to a measurable voltage signal through the Josephson phase. For representative benchmarks, this AB phase search reaches the minimum axion-photon coupling $g_{a\gamma \gamma }^{\mathrm{min}}\sim 7.8\times10^{-14}~\mathrm{GeV}^{-1}$ at axion mass $m_a\sim 10^{-10}~\mathrm{eV}$, with projected sensitivity that can improve on existing limits in that parameter space by roughly one to two orders of magnitude. We also identify a geometric phase observable in a Mach-Zehnder interferometer with an adiabatically rotating magnetic field, providing a proof-of-principle phase-based probe of meV-scale axions even when they do not constitute the dark matter, although sensitivity on the coupling remains weaker than current bounds with conservative tabletop benchmarks. Extending the analysis to a three level photon-axion quasiparticle (AQP)-axion system, with the AQP realized in a topological magnetic insulator, we find a potentially measurable THz Berry phase dominated by the AQP sector, furnishing a nontrivial validation of the formalism in a richer coupled system. These setups establish quantum phase observables as a useful new framework for axion searches, with immediate phenomenological promise in superconducting circuits and longer term potential in quantum enhanced interferometry.
\subsubsection*{Selected figures}
\begin{figure}[H]
\centering
\includegraphics[width=\linewidth,height=0.42\textheight,keepaspectratio]{/home/kf/research-assistant/data/cache/figures/http-arxiv.org-abs-2604.13181v1/figure\_1.png}
\end{figure}
\newpage
\subsection*{10. Targeted search for eccentric supermassive binary black holes in OJ 287 and nearby galaxy clusters with PPTA DR3}
\textbf{Focus:} galaxy clusters\\
\textbf{Authors:} Shi-Yi Zhao, Xingjiang Zhu, Jacob Cardinal Tremblay, Yiqin Chen, Małgorzata Curyło, Shi Dai, Valentina Di Marco, Pratyasha Gitika, George Hobbs, Simon C. -C. Ho, Xiao-Song Hu, Agastya Kapur, Wenhua Ling, Richard N. Manchester, Saurav Mishra, Daniel J. Reardon, Christopher J. Russell, Ryan M. Shannon, Sharon Mary Tomson, Jingbo Wang, Shuangqiang Wang, Andrew Zic\\
\textbf{Published:} 2026-04-14T18:00:13+00:00\\
\textbf{PDF:} \href{https://arxiv.org/pdf/2604.13173v1}{https://arxiv.org/pdf/2604.13173v1}\\
\textbf{Summary:}
We perform Bayesian targeted searches for continuous gravitational waves from eccentric supermassive binary black holes (SMBBHs) using the Parkes Pulsar Timing Array third data release (PPTA DR3). Six electromagnetically motivated sky directions are analyzed, including the blazar OJ\textasciitilde{}287 and five nearby galaxy clusters (Virgo, Fornax, Norma, Hercules, and Coma). No significant signals are found. For OJ 287, by explicitly incorporating orbital eccentricity (up to $e_0 = 0.8$) to robustly capture signal power spread across multiple harmonics, we constrain the total binary mass to $M_{\rm tot} < 5.25 \times 10^{10} M_{\odot}$ (95\textbackslash{}\% credible level). We also place upper limits on the chirp mass of potential SMBBHs residing in galaxy clusters. By combining these limits with independent black hole mass estimates, we place novel constraints on the allowed binary mass ratios for potential hosts such as M87 and NGC\textasciitilde{}4889. Specifically, our results exclude binaries with mass ratios $q \gtrsim 10^{-2}$ at around 10 nHz for these massive systems, effectively ruling out equal-mass black hole mergers in the sampled parameter space. These findings demonstrate the growing power of pulsar timing arrays to probe SMBBH populations.
\subsubsection*{Selected figures}
\begin{figure}[H]
\centering
\includegraphics[width=\linewidth,height=0.42\textheight,keepaspectratio]{/home/kf/research-assistant/data/cache/figures/http-arxiv.org-abs-2604.13173v1/figure\_1.png}
\end{figure}
\end{document}
